\section{Thesis Summary}
\label{sec:conclusion-summary}

This thesis has argued that causal credibility assessments are feasible in empirical software engineering through an explicit articulation of causal theory and assumptions.
The argument is developed through one tutorial chapter and three empirical chapters, each instantiating the four-step framework introduced in Chapter~\ref{chp:background}.
The contributions of each chapter are summarised below; the cross-chapter comparison in Section~\ref{sec:conclusion-comparison} makes their relative causal credibility explicit.

\subsection{A Tutorial on Causal Inference for Empirical Software Engineering}

Chapter~\ref{chp:background} provides the methodological infrastructure on which the three empirical chapters build.
A literature review of 5{,}341 main-track papers at ICSE, FSE, and ASE between 2015 and 2025 documents a structural gap in the field: 29\% of empirical studies are motivated by a causal research question, but only 1.9\% employ a recognised causal inference method, and the ratio has remained flat for a decade.
The tutorial closes this gap with a self-contained primer on the potential-outcomes framework, graphical causal models, and design-based identification strategies including difference-in-differences, instrumental variables, regression discontinuity, and panel fixed effects.
A worked example on the impact of AI coding tools illustrates how progressively stronger designs collapse a naive +911\% cross-sectional estimate to a credible +32--52\% difference-in-differences estimate, demonstrating that the choice of identification strategy is not a matter of statistical refinement but of which causal claim the data can actually license.

\subsection{Empirical Study I: Dependency Pinning and Supply Chain Security}

Chapter~\ref{chp:pinning} challenges the software engineering myth that pinning dependency versions improves project security against malicious package updates.
The chapter builds a counterfactual simulation that constructs both the pinning and the floating outcome for every project-time unit from the same \Code{package.json}, leveraging an experimental time-traveling feature of the npm dependency resolver to evaluate 10{,}000 npm packages and 10{,}000 GitHub repositories under both regimes.
A fixed-effects panel regression on the simulated data shows that pinning direct dependencies is futile or counter-productive in dependency graphs of 498 nodes or more, an effect driven by npm's deduplication algorithm and corroborated by a minimal worked example.
A complementary ecosystem-wide network propagation simulation shows that coordinated pinning at $\sim$100 important upstream packages reduces the spread of malicious updates by up to 75\%.
The methodological lesson is that identification by construction is achievable in observational software-engineering data when the underlying data-generating process can be simulated faithfully against a fixed historical snapshot.

\subsection{Empirical Study II: Fake GitHub Stars and Repository Promotion}

Chapter~\ref{chp:fake-stars} challenges the software engineering myth that buying fake GitHub stars provides sustainable promotional benefit.
The chapter introduces \SystemName, which mines GitHub's complete event metadata for two complementary signatures of inauthentic starring activity and identifies six million suspected fake stars across 26{,}254 repositories between 2019 and 2024.
A panel autoregression with repository fixed effects estimates the promotional and penalty effects of detected fake-star activity on subsequent organic star arrivals.
The analysis finds a short-run promotional bump of less than two months followed by a long-run penalty, with the majority of fake-star campaigns targeting short-lived phishing or malware repositories.
The methodological lesson is that even when the maintainer's treatment decision is endogenous and identification cannot be by construction, a panel-FE design paired with three targeted robustness checks---survivor stratification, topic stratification, and a lead-lag swap regression---can discipline the directions of residual biases and bound the policy-relevant claim, even though it cannot eliminate the underlying threats decisively.

\subsection{Empirical Study III: Cursor AI and Open-Source Project Outcomes}

Chapter~\ref{chp:cursor} challenges the software engineering myth that AI coding tools unambiguously improve developer productivity.
The chapter applies the \citet{borusyak2024revisiting} difference-in-differences estimator with propensity-score matching on six-month pre-adoption dynamics, comparing Cursor-adopting GitHub repositories with a matched control group of similar non-adopting repositories.
The DiD analysis finds a large but transient velocity gain in the first one to two months after adoption, alongside a substantial and persistent increase in static-analysis warnings and cognitive complexity.
A complementary panel generalised method-of-moments analysis decomposes the velocity-quality feedback and shows that the accumulated technical debt subsequently dampens future development velocity.
The methodological lesson is that a quasi-experimental design can reach close-to-strong causal credibility when the assignment mechanism cannot be controlled but matching, an estimator robust to heterogeneous treatment effects, and cross-direction GMM agreement substitute for the missing experimental variation.

\section{Comparing the Three Studies on Causal Credibility}
\label{sec:conclusion-comparison}

The three empirical chapters apply the same four-step framework but achieve distinctly different levels of causal credibility, governed by what the research settings and treatment assignment mechanisms permit and extended to the author's best rigour in their analyses.
The pinning study (Chapter~\ref{chp:pinning}) achieves identification by construction: both potential outcomes are simulated for every project from the same \Code{package.json}, so maintainer-choice confounders are blocked at the assignment mechanism rather than absorbed in expectation, and only simulation fidelity remains as an empirical argument.
The fake-stars study (Chapter~\ref{chp:fake-stars}) operates on fully observational panel data where the maintainer purchase decision is endogenous; three residual threats---reverse causality, topic-by-time confounding, and panel survivorship---are bounded by post-hoc robustness checks rather than by the design itself, and the identifying assumptions A1, A2, and A5 all require empirical argument.
The Cursor study (Chapter~\ref{chp:cursor}) sits between the two: It is stronger than the fake-stars study because the propensity-score matching procedure aligns observable pre-adoption trajectories, the \citet{borusyak2024revisiting} imputation estimator avoids forbidden comparisons under heterogeneous treatment effects, and the cross-direction GMM agreement substitutes for the missing experimental variation in distinguishing the velocity-quality feedback.
It is weaker than the pinning study because identification still rests on assumptions the design's construction does not enforce---parallel trends in the post-period and sequential exogeneity in the GMM---rather than on properties guaranteed by the assignment mechanism.

Table~\ref{tab:conclusion-comparison} summarises the comparison along the four dimensions of the framework.

\begin{table}[t]
    \scriptsize
    \centering
    \caption{Causal credibility comparison across the three empirical chapters.
    Each row reports how the chapter's treatment is assigned, the identifying assumptions, the diagnostic or robustness device that probes those assumptions, and the author's qualitative ranking of the resulting causal credibility.}
    \label{tab:conclusion-comparison}
    \begin{tabular}{p{2.1cm}p{3cm}p{7.4cm}p{1.6cm}}
    \toprule
    \textbf{Chapter} & \textbf{Dimension} & \textbf{Detail} & \textbf{Ranking} \\
    \midrule
    \textbf{Pinning (Ch.~\ref{chp:pinning})}
      & \textit{Treatment assignment} & Imposed deterministically by the simulator from each project's \Code{package.json} & {Strongest} \\
      & \textit{Identifying assumptions} & A1: simulation fidelity & \\
      & \textit{Diagnostic / robustness device} & Resolution-failure analysis on the 12--16\% of project-time units that fail to resolve; npm semantics validation against published resolver behavior & \\
    \addlinespace
    \textbf{Cursor (Ch.~\ref{chp:cursor})}
      & \textit{Treatment assignment} & Maintainer choice; matched on six-month pre-adoption trajectories & {Close to strong} \\
      & \textit{Identifying assumptions} & A1: parallel trends; B1: sequential exogeneity for the GMM & \\
      & \textit{Diagnostic / robustness device} & Pre-trend Wald test; Arellano-Bond AR(2) and Sargan tests; cross-estimator appendix (TWFE, \citet{callaway2021difference}) & \\
    \addlinespace
    \textbf{Fake stars (Ch.~\ref{chp:fake-stars})}
      & \textit{Treatment assignment} & Endogenous maintainer purchase decision, observed only through a noisy detection proxy & {Weakest, but informative} \\
      & \textit{Identifying assumptions} & A1: strict exogeneity; A2: no time-varying topic confounders; A5: panel completeness & \\
      & \textit{Diagnostic / robustness device} & R1: survivor-length stratification; R2: topic stratification; R3: lead-lag swap regression & \\
    \bottomrule
    \end{tabular}
\end{table}

The ordering in the rightmost column reflects a hierarchy that follows from the assignment mechanism: Identification by construction is preferable when the data-generating process can be simulated faithfully; design-based identification with targeted matching and an estimator robust to heterogeneous effects is the next-best option when the treatment is observed but not controlled; panel fixed effects with explicit robustness checks on the residual biases is the fallback when the treatment proxy is itself endogenous and noisy.
The ordering is not a verdict on the value of the contributions; all three chapters answer questions the field had not previously answered with quantitative evidence and important policy implications.

A surprise of the cross-study comparison is that policy survival is orthogonal to design strength.
The chapter with the weakest design (fake stars) yields a policy claim that survives most easily, because the residual biases on the short-run estimate point upward and would only strengthen the recommendation against buying fake stars if corrected.
The pinning chapter's policy claim survives because identification is strong by construction.
The Cursor chapter's policy claim survives because the most plausible residual confounders---a modernization push that bundles tool adoption with pace acceleration, mid-window adoption of competing AI tools, or pre-commit experimentation---all bias the estimates in directions that either reinforce the velocity-up and quality-down pattern or attenuate the velocity gains, leaving the qualitative recommendation intact.
Design strength and policy survival are therefore distinct dimensions, and downstream readers should resist the temptation to treat them as proxies for one another.

\section{Revisiting the Thesis Statement}
\label{sec:conclusion-thesis-statement}

\begin{quote}
\emph{Causal credibility assessments are feasible in empirical software engineering through an explicit articulation of causal theory and assumptions.}
\end{quote}

The four-step framework introduced in Chapter~\ref{chp:background} operationalizes this claim through a uniform procedure: 
\begin{enumerate}
  \item Derive an explicit causal theory from existing software engineering domain knowledge, typically as a directed acyclic graph.
  \item Define the causal estimand through the potential-outcomes framework and specify the assumptions under which it admits a causal interpretation.
  \item Honestly acknowledge and mitigate limitations and known caveats.
  \item Thoughtfully navigate alternative causal structures the data also admits.
\end{enumerate}
Each empirical chapter instantiates the four steps in a dedicated Causal Credibility Assessment section (Section~\ref{sec:pinning-credibility}, Section~\ref{sec:fake-stars-credibility}, Section~\ref{sec:cursor-credibility}).

The thesis statement claims feasibility, not certainty.
The fact that the same four-step framework operates uniformly across designs that differ as much as a counterfactual simulation, a staggered-adoption difference-in-differences, and a panel autoregression on observational data is the argument that the feasibility claim is broadly applicable rather than tied to any single design tradition.
I list a few important caveats regarding what this thesis does not claim and where the statement cannot be extended:
\begin{itemize}[leftmargin=20pt]
\item The framework does not eliminate causal uncertainty.
Each empirical chapter still has identifying assumptions that require empirical argument, and the credibility summaries name them honestly rather than rhetorically.
\item The framework does not apply to every empirical SE study.
The literature review in Chapter~\ref{chp:background} establishes that 29\% of empirical SE papers are related to a causal motivation; for the other 71\%, the causal credibility assessment framework is unnecessary.
\item The framework addresses identification, not measurement.
Persistent empirical controversies in SE---such as the programming-language and defect-proneness debate---remain unresolved partly because of identification failures (which this thesis addresses) and partly because of measurement challenges (which this thesis acknowledges but does not solve).
A causally clean estimate of an effect on a noisy proxy still leaves the underlying construct in dispute.
\item The framework does not replace existing data-driven empirical SE methods, nor does it imply that certain methods are inferior or superior to other methods.
This framework builds on top of the existing diverse realm of empirical methods in SE and adds an explicit statement of the causal theory, estimand, and assumptions that I believe will help transparent disclosure, iterative refinement, and cumulative knowledge building for a broad class of intervention-outcome problems prevalent in empirical software engineering research.
\item The framework prescribes transparent disclosure of the causal theory, not a particular formalism for it.
A DAG is one way to state a researcher's beliefs about the data-generating process, and different experts can legitimately hold different theories of the same phenomenon~\cite{DBLP:books/sp/08/SjobergDAH08, DBLP:journals/scp/StolF15}.
Some of the theories in this thesis---the velocity--quality feedback in Chapter~\ref{chp:cursor} and the momentum--purchase dynamics in Chapter~\ref{chp:fake-stars}---are temporal feedback systems that a static acyclic graph represents only after unrolling time, and for which cyclic structural causal models~\cite{bongers2021foundations}, time-varying treatment frameworks~\cite{robins1986new, hernan2020causal}, or system dynamics models~\cite{abdelhamid1991software, sterman2000business} may be the more natural formulation.
Section~\ref{sec:conclusion-transparency} extends this argument.
\end{itemize}

Within these limits, the three empirical chapters demonstrate that causal credibility assessment is feasible across substantively different software engineering questions, and that doing it explicitly produces more honest and more useful conclusions than the field's prevailing pattern of correlational claims accompanied by a perfunctory ``correlation does not imply causation'' caveat.

\section{Final Thoughts}
\label{sec:conclusion-final-thoughts}

\subsection{Revisiting Software Engineering Myths}

Chapter~\ref{chp:intro} framed the empirical software engineering field's recurring problem as one of \emph{myths}: widely held beliefs about best practices derived from intuition, experience, and conventional wisdom that have never been rigorously tested or, when tested, have produced inconclusive or contradictory results.
The three empirical chapters take three such myths in turn and apply the four-step framework to each.
The pinning myth---that pinning direct dependencies improves project security against malicious package updates---was endorsed by security tooling and authoritative scorecards (the OpenSSF Scorecard, the top-3 ranking in \citet{DBLP:conf/sp/LadisaPMB23}'s expert survey) and contested by a separate cohort of practitioners who view pinning as a maintenance liability~\cite{DBLP:journals/tse/JafariCAST22}; Chapter~\ref{chp:pinning} shows the security benefit reverses in large dependency graphs, complicating the myth in a way neither side of the prior debate had quantitative evidence for.
The fake-stars myth---that buying fake stars is an effective promotional shortcut---was held implicitly by the maintainers and startup founders who fund a real commercial market for fake stars; Chapter~\ref{chp:fake-stars} shows the promotional bump lasts less than two months and is followed by a long-run penalty, complicating the myth in a way that the practice's continued prevalence had not.
The AI-productivity myth---that AI coding tools unambiguously improve developer productivity---was reinforced by short-horizon controlled experiments and by vendor reports of multifold productivity gains; Chapter~\ref{chp:cursor} shows the velocity gain is transient and offset by sustained accumulation of technical debt, complicating the myth in a way that single-task experiments could not detect.

In each case, the myth was not simply true or simply false; it was a coarse claim about a complicated mechanism that admitted a partial reading along one dimension and a contrary reading along another.
The four-step framework's contribution is therefore not to declare each myth dispelled but to specify the conditions under which the myth holds, the conditions under which it does not, and the identifying assumptions on which both directions of the claim rest.
\emph{This conditional reframing is more useful for the field than dispelling the myths outright.}
Practitioners and researchers can then update their beliefs along the specific dimension the evidence speaks to rather than wholesale, and downstream readers can engage with the qualified claim without having to choose between accepting or rejecting the entire study.
% This reframing also explains why the field's prevailing pattern of correlational findings with a perfunctory ``correlation does not imply causation'' caveat has done little to dispel myths over decades: An unconditional finding cannot displace an unconditional belief, and only a conditional, identification-aware claim can engage with a myth at the resolution at which the myth is actually held.

\subsection{Revisiting the Downstream Misuse Problem}

A central worry that motivates this thesis is that even carefully hedged empirical results are routinely read by downstream readers as causal claims, and that the field lacks a shared framework for evaluating whether such a reading is warranted.
The Ray et al.\ controversy documented in Chapter~\ref{chp:intro} illustrates the pattern: the original authors framed their findings as associational, but downstream practitioners and follow-up papers cited the work as if it had established a causal effect of programming language on defect proneness.
The four-step framework cannot prevent miscitation, but it offers a partial response along three lines.
For producers of empirical work, the framework imposes a discipline that forces the identifying assumptions into the open before the study reaches downstream readers.
For consumers of empirical work, the framework provides a checklist for distinguishing studies whose authors have stress-tested their identification and reported what survived from studies that have not.
For misciting downstream readers, the framework cannot police behaviour, but it makes the gap between what a study claims and what a downstream reader infers visible and challengeable.

The fake-stars chapter is the clearest illustration of this asymmetry.
A downstream reader who reads only the abstract still receives the policy claim that buying fake stars is ineffective; a downstream reader who engages with Section~\ref{sec:fake-stars-credibility} understands why the claim survives, what it depends on, and what residual threats a careful reader would still want to inspect.
The framework is therefore a piece of communicative infrastructure as much as it is a methodological one: it gives a downstream reader the tools to read a study honestly even when that reader is not in a position to redo the analysis.

\subsection{Assumptions Are Universal, and Transparent Disclosure Is the Key}
\label{sec:conclusion-transparency}

Every causal inference rests on assumptions, and observational designs are not uniquely vulnerable in this respect.
Randomized controlled trials rest on assumptions of stable unit treatment values, full compliance, no spillover between units, construct validity of the chosen outcome measure, and generalisability beyond the trial population.
The novelty of the four-step framework is not that it makes assumptions visible---every method does, somewhere---but that it makes them visible uniformly across the observational designs that the empirical SE field has historically treated as if they had no assumptions to articulate.
The shift the thesis advocates is therefore one of expectation rather than of method: A study that claims a causal effect should be expected to name the assumptions the claim rests on, just as a randomised trial is expected to report its compliance rate.
Still, this shift will not satisfy the most cynical critic, who can refuse to believe any observational study, and it will not protect against the uninformed reader, who can misread any descriptive result as causal.
Neither failure mode is the framework's responsibility to solve.
The framework's value is for the practitioner in between---the researcher, reviewer, or practitioner-engineer who is willing to engage with the evidence but needs a structured way to calibrate confidence in what the evidence supports.

The same expectation of disclosure applies to the causal theory itself, and here the framework's commitment is to transparency rather than to any single representation.
Chapter~\ref{chp:background} recommends DAGs because they make conditioning decisions inspectable and because the back-door criterion turns a drawn graph into a covariate-selection rule; the three empirical chapters follow this recommendation.
A DAG is nevertheless one formalism among several for encoding what a researcher believes about the data-generating process.
The potential-outcomes tradition in economics states the same beliefs as assumptions about the assignment mechanism and rarely draws a graph~\cite{imbens2020potential}; structural econometrics states them as systems of equations with explicit functional form~\cite{heckman2015causal}; and the medical-statistics literature has questioned whether an acyclic graph is the right abstraction for mechanisms that unfold continuously in time~\cite{aalen2016can}.
The distinction matters for software engineering because many of the field's most interesting theories are not completely acyclic.
The Cursor chapter's central finding is a feedback loop---velocity affects quality within a period and quality affects velocity in the next---and the fake-stars chapter's principal residual threat is a reverse-causal loop between organic momentum and purchase decisions.
Both chapters handle these loops by unrolling time so that the graph remains acyclic, which is the standard treatment of time-varying exposures~\cite{robins1986new, hernan2020causal}, but a researcher whose theory is fundamentally about equilibria, accumulation, or delayed feedback may find a cyclic structural causal model~\cite{bongers2021foundations} or a system dynamics model in the tradition of \citet{abdelhamid1991software} and \citet{sterman2000business} a more faithful statement of that theory.
This thesis leaves these as potential future work.

Different experts also hold different theories of the same phenomenon, and the software engineering literature has long recognized that the field's theories are plural, mostly implicit, and unevenly used.
\citet{DBLP:journals/tse/HannaySD07} found that fewer than a quarter of the software engineering experiments published over a decade used any theory to explain the cause--effect relationship under investigation, and subsequent calls for theory-oriented software engineering~\cite{DBLP:journals/software/JohnsonEJ12, DBLP:journals/scp/StolF15, DBLP:journals/tse/Ralph19} have argued that the discipline needs a shared vocabulary for proposing, comparing, and refining such theories~\cite{DBLP:books/sp/08/SjobergDAH08}.
Causal credibility assessment supplies one concrete use for that vocabulary: Whatever formalism a researcher chooses, the theory must be stated explicitly enough that a second expert holding a different theory can see where the two disagree, and the fourth step of the framework is the place where those disagreements are meant to surface.
The requirement is therefore not that every study draw a DAG, but that every study make its theory a disclosed and contestable object rather than an implicit input to a model specification.

Transparent disclosure has a further benefit that the preceding discussion has only implied: It is the precondition for iterative refinement.
An assumption that is named can be targeted by a later design; a theory that is stated can be contested by a second expert; and a limitation whose direction of bias is recorded tells the next study which way its estimate needs to move.
Each generation of empirical SE work that engages explicitly with causal credibility therefore hands the next generation both a stronger toolkit and a sharper list of what remains to be shown; the field's path forward is iterative refinement, not a single final resolution.

\subsection{On Practitioners' Beliefs and the Impact of Empirical SE Research}
\label{sec:conclusion-practitioner-beliefs}

A reasonable objection to the three empirical chapters is that their estimates match nobody's prior.
This is because practitioners form their beliefs primarily from personal experience rather than from published evidence~\cite{DBLP:conf/icse/Devanbu0B16, DBLP:conf/esem/PassosBCM11}, they are persuaded more readily by local opinion than by empirical findings~\cite{DBLP:conf/isese/RainerHB03}, and their beliefs are not consistently supported by data from the very projects that formed them~\cite{DBLP:journals/ese/ShrikanthNFM21}.
On the other hand, the pinning chapter's estimate---a security benefit that reverses in dependency graphs of 498 nodes or more---is a population-level, conditional quantity that neither the security cohort nor the maintenance cohort in the prior debate held a belief about, and the same is true of the two-month promotional horizon in the fake-stars chapter and the transient velocity gain in the Cursor chapter.
If the estimate corresponds to no belief a practitioner actually holds, what is the study for?

This thesis takes the position that empirical software engineering research is \emph{supplementary} to practitioner judgment, and that its best achievable impact is to shift practitioners' beliefs by the amount that well-informed and well-executed evidence warrants.
The position has a precedent.
Evidence-based medicine defined itself as the integration of individual clinical expertise with the best available external evidence and explicitly rejected the reading that evidence replaces judgment~\cite{sackett1996evidence}; the evidence-based software engineering program imported the same stance~\cite{DBLP:conf/icse/KitchenhamDJ04, DBLP:journals/software/DybaKJ05}.
Research is supplementary not because it is weaker than experience but because the two sample the data-generating process differently.
A practitioner's experience samples a narrow slice---their own projects, ecosystem, and era---at high resolution, and it is frequently accurate within that slice.
Empirical research samples broadly across projects at lower resolution and recovers population-level and conditional structure that no individual's experience can reach.
For example, the pinning debate illustrates the complementarity: The security cohort and the maintenance cohort each generalized from the projects they had seen, both were right within their slice, and the dependency-graph-size condition that reconciles them is visible only from an ecosystem-scale sample.
That the estimate matches no prior is therefore not a defect of the study; it is the information the study adds.
A practitioner who combines the two sources arrives at a belief that neither source could have produced alone, and the impact of a study is the size of that update rather than the estimate it reports.

The Bayesian framing introduced in Chapter~\ref{chp:intro} makes the position precise and clarifies what causal credibility contributes to it.
A practitioner who assigns a low prior probability to a claim---which is what holding a myth means---updates toward the claim only when the observed evidence is substantially more probable under the claim than under its alternatives~\cite{howson2006scientific}; this is the sense in which extraordinary claims require extraordinary evidence.
The consequence that matters for this thesis is that the likelihood ratio is a property of the research design, not of the point estimate.
If an observed effect is nearly as probable under confounding, reverse causality, or selection as under the causal claim, the likelihood ratio is close to one, and the estimate should not move a calibrated practitioner however large it is.
This is exactly the discounting Chapter~\ref{chp:intro} describes: When a controversial finding fails to rule out alternative explanations, practitioners become skeptical of it altogether, and that skepticism is warranted.
Steps 3 and 4 of the framework---bounding residual biases and navigating alternative causal structures---are where alternatives are made improbable, so a causal credibility assessment is, in Bayesian terms, an honest report of the likelihood ratio the study is entitled to claim.
Extraordinary evidence is built by design rather than by effect size, and iterative refinement compounds it: Each subsequent study that targets an assumption the previous one left open removes an alternative under which the earlier evidence would have been equally likely, so the accumulated likelihood ratio across a research program grows in a way that a sequence of correlational studies restating the same association under the same unaddressed alternatives cannot.
A myth held with confidence will therefore not be displaced by a single study, and the framework does not promise that it will; it promises that each study's contribution to the eventual displacement is stated and can be built upon.

I make three additional clarifications to keep this position from overreaching:
\begin{itemize}
  \item First, the goal is calibration, not persuasion in a predetermined direction.
A research program whose only acceptable output is a dispelled myth would inflate its claims, and a practitioner population that has learned to discount inflated claims would assign every subsequent study a likelihood ratio near one, corroding the channel the program depends on.
The framework must remain equally willing to report that a myth survives.
  \item Second, practitioners hold heterogeneous priors, so there is no single threshold of extraordinary evidence; the threshold is set by the most skeptical reasonable prior, and the conditional claims discussed above are what let a study address each prior at its own resolution.
For a practitioner whose belief concerns a condition the study does not analyze, the researcher must provide the honest answer that the study does not speak to them and disclose the boundries transparently.
  \item Third, the belief update rarely happens through the paper itself.
Practitioners rate much software engineering research as relevant~\cite{DBLP:conf/sigsoft/LoNZ15}, yet published findings remain a minor source of their beliefs relative to experience and peers~\cite{DBLP:conf/icse/Devanbu0B16, DBLP:journals/ese/GarousiBO20}.
Beliefs travel through intermediaries---security scorecards, tool defaults, platform policies, educators, and colleagues---and the audience whose beliefs a study most needs to move is often relevant practitioners (e.g., maintainers of the OpenSSF Scorecard for the pinning study and senior managers deciding tool adoption for the Cursor study) rather than any individual developer.
The communicative role of the framework discussed above matters most for these intermediaries.
\end{itemize}

\subsection{Closing Posture: How Should the Field Move Forward?}
\label{sec:conclusion-closing}

The thesis argues that causal credibility assessment is feasible; whether the field moves as a result is a separate question, and there are two symmetric ways it could go wrong.
The first failure mode is inertia: The field takes note of the framework and changes nothing, and the 1.9\% share of causal-method papers documented in Chapter~\ref{chp:background} stays where it has stood for a decade.
The second failure mode is over-correction: Causal inference becomes a fashion, peer review develops an \emph{identification police} that rejects observational studies for lacking an instrument, a discontinuity, or an experiment, and the field loses the descriptive, exploratory, and qualitative work on which causal theories depend.
Both failure modes have precedents in fields that adopted causal inference earlier than software engineering, and their experience shapes my position here.

The identification-police concern is the one often raised against this thesis, and the economics profession's experience supplies the answer.
The credibility revolution in empirical economics~\cite{angrist2010credibility} provoked a sustained critique that design-based methods answer narrow questions with local estimands, privilege identification over substantive theory, and reduce methodology to a hierarchy in which only experiments and instruments count~\cite{deaton2010instruments, heckman2010comparing, deaton2018understanding}.
The productive part of that exchange was the discipline the debate imposed on both sides: A study must state what its design identifies, for which population, and under which assumptions, and a critic must state which assumption fails and in which direction~\cite{imbens2010better}.
The unproductive part was the reflex of dismissing a study for the design it did not use rather than engaging with the assumptions it did.
None of the three empirical chapters in this thesis uses an instrument or a discontinuity, and a reviewer who applied the hierarchy as a rule would reject all three.
Each chapter also names its assumptions, so the valid criticism of each is specific: that strict exogeneity in Chapter~\ref{chp:fake-stars} fails because momentum drives purchases in a way the lead-lag swap does not bound, that parallel trends in Chapter~\ref{chp:cursor} fails because a modernization push coincides with adoption, or that the simulation in Chapter~\ref{chp:pinning} misrepresents resolver behavior for some class of projects.
In all cases, the framework cannot disarm criticism but may rather give weapons through the exact transparent disclosure process I advocate throughout the thesis.

This is why an empirical standard is the right institutional response.
The ACM SIGSOFT Empirical Standards~\cite{DBLP:journals/corr/abs-2010-03525} address reviewers as well as authors, and each standard enumerates acceptable deviations and invalid criticisms alongside essential attributes.
The standard for causal inquiries in Appendix~\ref{sec:appendix-standard} inherits that design: It names the absence of a randomized experiment, the existence of unmeasured confounders, a local estimand, a hedged conclusion, and the use of an established rather than a novel estimator as grounds on which a reviewer may not reject.
Importantly, the standard makes it explicit upfront that it does not apply to descriptive, exploratory, and qualitative work that does not seek to quantitatively investigate intervention-outcome questions. 
I should note that the field also already has a police of the opposite kind.
The prevailing reflex that treats any causal language in an observational study as a defect, and that is satisfied by a perfunctory ``correlation does not imply causation'' caveat, is an \emph{association police} whose effect epidemiology has documented: Euphemism hides the causal question the study was asking and does nothing for the inference~\cite{hernan2018cword}.
For both kind of ``police'' in peer reviews, an empirical standard that clearly outlines invalid criticisms is the answer.

The inertia concern is the more serious of the two, and the thesis's own evidence supports taking it seriously.
The field of medicine offers the sharper warning: Evidence-based medicine has had a mature evidence base and formal guidelines for three decades~\cite{sackett1996evidence}, and a substantial literature documents that well-evidenced guidelines do not change practice wholesale.
\citet{cabana1999why} classify the barriers as lack of awareness and familiarity, lack of agreement, lack of self-efficacy, lack of outcome expectancy, inertia of previous practice, and external constraints, and \citet{grol2003best} show that passive dissemination of evidence is largely ineffective on its own and that durable change requires interventions targeted at the specific barriers of a specific setting.
Read against this classification, the question of what the field can do beyond adopting methods has a concrete answer: Target the barriers one at a time.
The tutorial chapter targets awareness and familiarity, and the worked example, the FAQ in Appendix~\ref{sec:appendix-faq}, and the three empirical chapters as reusable templates target self-efficacy.
The three chapters also speak to outcome expectancy by producing findings that no correlational design could have delivered: an effect that reverses with dependency-graph size, a promotional bump that turns into a penalty, and a velocity gain that technical debt subsequently consumes.
The empirical standard is aimed at the external constraint that matters most to researchers, which is what reviewers and venues expect.
Inertia is the barrier that none of these moves individually, and it yields only when outcome expectancy and venue norms move together.\footnote{
  I hold reasonable optimism as adoption is already visible.
  Within eight months of the Cursor study's preprint, several independent groups reused its design: \citet{DBLP:journals/corr/abs-2606-13298} apply the same staggered difference-in-differences design and \citet{borusyak2024revisiting} imputation estimator to architectural smells and describe the Cursor study as their closest predecessor; \citet{DBLP:journals/corr/abs-2607-01810} use the same imputation estimator and the same cognitive-complexity metric to test whether the complexity increase crowds out newcomers; and \citet{DBLP:journals/corr/abs-2606-26289} adopt its propensity-score matching procedure to study contributor ecosystems after agent adoption.
  \citet{DBLP:journals/corr/abs-2608-03329} apply matching and difference-in-differences to the effect of AI governance policies in open-source projects.
  Independent of this thesis, causal designs are also entering the field's methodological discourse, including a registered report on the causal effect of code coverage~\cite{DBLP:journals/corr/abs-2602-03585}, a structural-causal-model framework for empirical SE~\cite{DBLP:journals/corr/abs-2605-28482}, and a tutorial on applied causal inference in requirements engineering~\cite{DBLP:journals/corr/abs-2511-03875}.
  These are preprints that have not yet completed peer review, and a handful of studies does not move a decade-long adoption rate; they nevertheless indicate that the outcome-expectancy barrier has started to give.
}

Ultimately, this thesis does not claim to resolve the empirical software engineering field's longstanding disconnect between causal ambition and methodological practice.
It claims that the disconnect is bridgeable through an explicit articulation of causal theory and assumptions, and that the four-step framework is one usable operationalisation of that bridge.
The three empirical chapters demonstrate that the framework is feasible across substantively different questions; the cross-chapter comparison demonstrates that feasibility does not collapse credibility differences across designs but makes them inspectable.
If future empirical SE research adopts this kind of explicit credibility assessment, even partially, the field will become better equipped to do what it has historically struggled to do: engage the myths that practitioners hold at the resolution at which they are held, by accumulating transparently justified causal evidence with a clear statement of what the evidence supports, under which assumptions, and to what extent.

